\section{Introduction}
\label{sec:intro}

The problem of Course Timetabling in a university is the problem of assigning many classes with a set of rooms, teachers and time constraints to a weekly timetable in a grid, under many different hard and soft constraints. The problem is NP-hard, which means that we do not know any polynomial time exact solution or have practical systems that use heuristic or metaheuristic search~\cite{schaerf1999survey,muller2007}.

COEP Technological University offers a challenging example of this. The Department of Computer Engineering has four undergraduate years and a postgraduate programme that each consist of lecture divisions and then laboratory batches, with departmental electives, minor degree courses offered by other disciplines, open electives and reserved first-year batches. Laboratory sections are parallel courses, courses are parallel lectures, and a scheduling unit is not the course itself, but the section of it, to the extent that the scheduling unit is fit to be used, or not, depends on the relations between the sections, not the course itself.

\subsection{The school values that were set by Semester~VI}

The previous step in this project was to assess UniTime, an academic scheduling system developed as an open source, web-based platform using the CPSolver constraint library, for suitability as a platform. It installed the system in its own environment and loaded representative test data into it both via the web interface and via the XML import feature, and translated the institution's scheduling rules into UniTime's constraint model. The conclusion was that all the constraints found by the team were representable in that model and none of the constraints needed to be relaxed or post-processed externally.

Upon completion of Semester~VI, there was no institutional data, no real solution found from the course structure, and no measure for the quality of the solution. All the data pipeline, student records, and the instructor-facing interface were all described in the future tense.

\subsection{There are 1.2 credits earned for Semester~VII}

The current phase is the \emph{taasika-v3}, which addresses that. The entire work is organized in five contributions, each of which is reported along with the evidence that supports it:

\begin{enumerate}
  \item \textbf{A reproducible Taasika\,$\rightarrow$\,UniTime data pipeline.} Eight Python
    generators read a snapshot of the university's existing Taasika scheduling database and emit
    the sixteen ordered XML documents that UniTime's Data Exchange interface consumes, followed by
    three validators and an archive step, behind a single command
    (Section~\ref{sec:pipeline}).
  \item \textbf{Institution-specific constraints encoded as data.} Minor-degree and open-elective
    block times, division and batch structure, elective spread, and dynamic lunch breaks for
    students and instructors are generated into the preference file and solver configuration rather
    than entered by hand (Section~\ref{sec:constraints}).
  \item \textbf{An unbiased assessment system.} Because a solver cannot be audited with its own
    counters (Section~\ref{sec:evaluation}), a 1{,}044-line evaluator is used to score any exported
    timetable against the same input XML with 10 hard checks, 3 preference checks and a set of
    efficiency metrics.
  \item \textbf{The accurate description of the solution quality.} A number of configurations
    experienced no student conflict reported, all of which were found to be artefacts. For the last
    30 weeks, student enrolments have been fixed, and double-bookings of instructors and rooms have
    been avoided: 290 of 291 classes have been scheduled in a way that avoids instructor or room
    double-booking; 40 of 1{,}130 students are in conflict (Section~\ref{sec:experiments}).
  \item \textbf{Reinforcement learning for the solver search.} A policy is a choice from eight
    neighbourhood operators for the CPSolver, which are implemented within the search loop. The
    layer is implemented and instrumented, and is the result of a negative answer to the research
    question (Section~\ref{sec:rl}).
\end{enumerate}

The project is referred to throughout as \emph{taasika-v3}. The source repository retains an earlier
name, \texttt{coep-unitime-pilot}, which should be read as a historical artefact rather than as a
description of the current scope.
